\documentclass[a4paper,twoside,12pt,AutoFakeBold]{ctexart}
\newcommand\specialsectioning{\setcounter{secnumdepth}{-2}}
\specialsectioning
\usepackage[center]{titlesec}
\usepackage{indentfirst}
\setlength{\parindent}{2em}
\usepackage{fancyhdr}
\pagestyle{fancy}%fancy style
\usepackage{caption}
\fancyhf{}%清空页眉页脚
\fancyhead[LE,RO]{\thepage}%页码位置：偶数页居左，奇数页居右
\fancyfoot[RO,RE]{\textit{Tongji's seminar of Marxism}}% 设置页脚：在每页的右下脚以斜体显示书名
\usepackage{graphicx}
\setlength{\headheight}{15pt}%解决页眉warnings
\usepackage{amsmath}
\renewcommand{\headrulewidth}{0pt} % 页眉与正文之间的水平线粗细
\renewcommand{\footrulewidth}{0pt}
\usepackage{tcolorbox}%文本框宏包
\usepackage{changepage}%设置引用段落左右侧缩进
\usepackage{nameref}
\usepackage{tabularx}
\usepackage[colorlinks,linkcolor=blue,
anchorcolor=green,citecolor=green]{hyperref}%设置超链接
\usepackage{float}
\title{读书会记录}
\author{\kaishu 同济大学马克思主义经典著作研讨会\\
%\kaishu 指导教师\\
\kaishu 吴云飞~编}
\date{2026年秋季学期}
\usepackage{perpage}
\MakePerPage{footnote}
\usepackage{soul}%下划线宏包
\usepackage{eso-pic} %添加背景图片
\newcommand\BackgroundPic{%
\put(-20,-90){ % 调整图片的位置
\parbox[b][\paperheight]{\paperwidth}{%
\vfill
\centering
\includegraphics[width=4in,angle=0, % 设置图片宽度，旋转0度
keepaspectratio]{badge.png}
\vfill
}}}


\begin{document}
\maketitle
\AddToShipoutPicture*{\BackgroundPic}  % 在当前页面添加背景图片
\newpage


\tableofcontents%目录

\newpage

\section{序言}

\begin{adjustwidth}{2em}{2em}
\qquad\fangsong 
本书是同济大学马克思主义经典著作研究与讨论会的一个内部记录，诸多观点与看法可能会欠缺专业性，因而这是一个非教学性质的内容记录。本书全部内容均来自我们的研讨会中的发言与讨论以及后续编者所进行的思考，如有雷同，纯属巧合，本书最终解释权归本研讨会与本书编者\footnote{编者注：我的邮箱是\href{mailto@tongji.edu.cn}{wuyunfei2000@tongji.edu.cn}，如果您在阅读本书之后有相应的意见与建议，抑或是想进行一些学术方面的交流与探讨，请通过该邮箱与我联系。}所有。

本书采用\LaTeX{}编写，\TeX{}源代码已开源，且暂由编者吴云飞托管至其本人的Github仓库之中，供各位成员获取:\url{https://github.com/Imheaven233/Tongji-University-Seminar-on-Marxism}。

\end{adjustwidth}

\newpage

\section{第一期：《<黑格尔法哲学批判>导言》（1）}

\subsection{学习提示}\label{sec:1}
《<黑格尔法哲学批判>导言》是马克思从唯心主义向唯物主义、从革命民主主义向共产主义转变过程中的重要著作。该文章撰写于1843年10月中至12月中，并于1844年2月在《德法年检》公开发表。\footnote{编者注：《导言》主要写于巴黎，而《黑格尔法哲学批判》主体手稿主要写于克罗茨纳赫。}

这篇文章不长，接下来我们会对这篇文章的完整全文进行讨论。

在本次讨论会中，您将会了解到（或带着以下问题去阅读）：

\textbf{1.}为什么马克思要强调宗教批判的重要性？

\textbf{2.}青年黑格尔派（起源于哲学的理论政治派）对宗教进行的批判为何存在着局限？

\textbf{3.}马克思的观点是什么？

\textbf{4.}为什么有学者认为《导言》是马克思从革命民主主义向共产主义、从唯心主义向唯物主义转变过程中的代表性文献？这种观点有道理吗？

\subsection{读书会记录}
2026年9月19日13：30—15：00，我们在同济大学马克思主义学院413教室开展了第一期的读书会。本期读书会我们对《<黑格尔法哲学批判>导言》进行了阅读\footnote{编者注：没有完整读完，剩下的内容会安排在下期继续讨论。}，并展开了充分的讨论。具体讨论内容可概括如下：



\subsubsection{一、如何理解“对宗教的批判是其他一切批判的前提”？}\label{1.1}
首先，马克思在在开头段指出：
\begin{adjustwidth}{2em}{2em}
\qquad\fangsong
“就德国来说，对宗教的批判基本上已经结束；而对宗教的批判是其他一切批判的前提。”
\end{adjustwidth}

马克思为什么要说“对宗教的批判是其他一切批判的前提”呢？对此我们进行了讨论。

我们认为，这首先与马克思对宗教和现实社会关系的理解有关。在马克思看来，宗教并不是凭空产生的，而是现实社会的产物。宗教之所以表现为一种“颠倒的社会意识”，根本原因在于产生宗教的现实世界本身就是颠倒的。因此，对宗教的批判虽然直接表现为对颠倒的社会意识的批判，但其意义并不止于此，而在于通过对这种颠倒的社会意识的批判，进一步追问产生这种意识的现实社会基础。也就是说，对宗教的批判构成了从意识批判走向现实批判的一个起点。

同时，我们认为这句话也需要放到当时德国的现实背景中来理解。19世纪三四十年代，宗教在德国社会生活和思想领域中仍然具有十分显著的影响，它不仅是一种普遍存在的社会意识形式，而且在相当程度上承载和反映着旧社会、旧制度的思想因素。也正因如此，宗教问题并不是一个孤立的观念问题，而是当时德国现实社会矛盾在意识领域中的集中表现。因此，要批判当时德国的那种旧的社会现实，就必须首先从其最为显著的意识形态表现入手，即对宗教进行批判。换言之，对宗教的批判之所以成为“其他一切批判的前提”，不仅因为宗教是一种颠倒的社会意识，更因为在当时德国，宗教与旧社会秩序之间存在着密切联系，对宗教的批判因而构成了进一步批判法、政治乃至整个现实社会的重要起点。

正因如此，马克思随后才会提出：

\begin{adjustwidth}{2em}{2em}
\qquad\fangsong
“因此，真理的彼岸世界消逝以后，历史的任务就是确立此岸世界的真理。人的自我异化的神圣形象被揭穿以后，揭露具有非神圣形象的自我异化，就成了为历史服务的哲学的迫切任务。于是，对天国的批判变成对尘世的批判，对宗教的批判变成对法的批判，对神学的批判变成对政治的批判。”
\end{adjustwidth}

我们不难看出，马克思这里所说的“人的自我异化的神圣形象”其实指的就是宗教，揭穿“人的自我异化的神圣形象”，指的就是对宗教的批判。因此，马克思会说，揭穿这种“神圣形象”之后，还需要进一步揭露“非神圣形象的自我异化”，也就是还需要进一步批判产生这种“自我异化神圣形象（宗教）”的社会现实。

\subsubsection{二、马克思是怎么看待“人的本质”的？}

马克思在《导言》中概括了当时德国哲学界对宗教批判的基本前提：

\begin{adjustwidth}{2em}{2em}
\qquad\fangsong
“反宗教的批判的根据是：人创造了宗教，而不是宗教创造人。就是说，宗教是还没有获得自身或已经再度丧失自身的人的自我意识和自我感觉。”
\end{adjustwidth}

这句话的意思是很好理解的。马克思想表明，当时的德国哲学界对宗教批判遵循的基本逻辑是：人创造了宗教，而不是宗教创造了人。循此逻辑继续推演下去，就是说，宗教是“人”（人的本质）的反映，“人”不能被宗教束缚住，因此批判宗教的目的就是让宗教真正的复归于“人”（人的本质）。

乍一看，我们可能觉得这种逻辑没啥问题。但我们在后文能够看到，马克思对这种观点显然是不满意的。马克思紧接着论述到：

\begin{adjustwidth}{2em}{2em}
\qquad\fangsong
“但是，人不是抽象的蛰居于世界之外的存在物。人就是人的世界，就是国家，社会。这个国家、这个社会产生了宗教，一种颠倒的世界意识，因为它们就是颠倒的世界……宗教是人的本质在幻想中的实现，因为人的本质不具有真正的现实性。”
\end{adjustwidth}

值得注意的是，马克思在这里使用的是“但是”，体现了一种转折。前面说“人创造了宗教”，马克思对此并没有否定。但之后我们能够感觉到，马克思关注的真正的问题在于，这里的“人”究竟应当怎样理解。如果只是从抽象概念的角度去理解“人”，进而把宗教从“神”重新归结为这样的一个抽象的“人”，那么在马克思看来这种批判仍然没有真正触及宗教产生的现实根源（本质维度）。

因此马克思会马上强调：“人不是抽象的蛰居于世界之外的存在物。人就是人的世界，就是国家，社会。”我们可以看到，马克思在这里并没有像其他哲学家那样对所谓“人”或“人的本质”作出定义，马克思只是说“人不是抽象的存在物，人就是人的世界，就是国家，社会。”在这里，不难看出，马克思对“人”或“人的本质”的理解并不是从哲学、形而上学的角度去理解的，而是从实实在在的感性现实中去理解的。因此，马克思会着重强调：“人的本质不具有真正的现实性。”

关于这一点，我们认为可以联系马克思的《关于费尔巴哈的提纲》第六条中的内容加深理解。在《提纲》的第六条，马克思是这么说的：

\begin{adjustwidth}{2em}{2em}
\qquad\fangsong
“费尔巴哈把宗教的本质归结于人的本质。但是，人的本质不是单个人所固有的抽象物，\textbf{在其现实性上}，它是一切社会关系的总和。”
\end{adjustwidth}

乍一看，《提纲》第六条中的论述似乎和《导言》中的叙述存在着矛盾。在《导言》中，马克思说“人的本质不具有真正的现实性”。而到了《提纲》第六条，马克思却说：“人的本质在现实性上是一切社会关系的总和。”那么，在马克思看来，“人的本质”到底有没有现实性？马克思是不是前后矛盾了？

实际上，马克思非但没有前后矛盾，反而这两处的论述实际上表达的是同一回事。根据我们前面的讨论，我们可以得到一个信息，马克思不是从哲学、形而上学的角度去谈论所谓“人”或“人的本质”的，而是从感性现实中对其进行把握。因此，马克思之所以会在《导言》中强调“人的本质不具有真正的现实性”，是想表明，实际上根本不存在一种所谓的“人的本质”的形而上学的抽象概念，因而这种“抽象的概念”是“非现实性”的。

当我们理解了这一点之后，我们再回到《提纲》第六条中的论述。马克思的前半句说“人的本质不是单个人所固有的抽象物”，就能够反映出我们在上面讨论中得到的看法，即不存在一种所谓的“人的本质”的形而上学的抽象概念。而后半句，马克思又说“\textbf{在其现实性上}，它是一切社会关系的总和。”事实上，这里的“现实性”指的并不是作为抽象概念的“人的本质”的现实性，而是指具体层面的人在感性现实中的具体呈现。

最后，用通俗易懂的话来总结一下马克思关于“人的本质”的看法：
\textbf{“人的本质”是形而上学语境下的抽象物，因而是非现实性的存在，但这种“非现实性的存在”却是在“现实性”层面中产生的，而产生这种“非现实性”的“现实性”就是一切社会关系的总和。}

{\kaishu 因此，如果有人问马克思：“究竟什么是人的本质？”马克思会回答他：“我也不知道，因为压根不存在人的本质这种抽象物。”如果那个人不满意马克思的回答，继续追问：“不行，你的回答我不满意，无论如何你都得对人的本质下个定义。”马克思就会说：“我没法直接给人的本质下定义，因为它是不存在的抽象概念，但如果你非得让我对这个抽象概念作出定义，我只能告诉你，在现实性上，它表现为一切社会关系的总和。”}


\subsubsection{三、德国革命的特殊性}

在“\nameref{1.1}”中，我们讨论了马克思对待宗教批判的基本态度，即宗教批判是有必要的，但不能仅仅停留于宗教批判，还必须进一步转向现实批判。接下来，马克思讨论了德国的具体情况。马克思首先说：

\begin{adjustwidth}{2em}{2em}
\qquad\fangsong
“随导言之后将要作的探讨——这是为这项工作尽的一份力——首先不是联系原本，而是联系副本即联系德国的国家哲学和法哲学来进行的。其所以如此，正是因为这一探讨是联系德国进行的。”
\end{adjustwidth}

我们认为，这里的“原本”主要指现代国家现实存在的政治和社会关系，而“副本”则是指德国关于现代国家的国家哲学和法哲学。就当时欧洲的发展状况而言，英法等国已经较早进入现代资本主义社会，现代国家的政治制度和社会关系已经在现实中得到较为充分的发展；德国则仍然保留着大量封建专制因素，其现实政治发展明显落后于这些现代国家。从这个意义上说，德国并不具备与英法相同的现实批判对象。

但是，这里随之产生了一个问题：既然马克思前面已经强调，批判应当由宗教批判转向现实批判，那么为什么到了德国，却又提出要首先批判德国的国家哲学和法哲学呢？

事实上，这恰恰反映了德国发展的特殊性。德国虽然在现实政治发展上落后于英法等国，但在哲学上却并没有完全落后于时代。尤其是德国哲学已经在思想中对现代国家、政治和法的问题进行了高度抽象和系统化的表达。也就是说，德国人在现实中尚未充分经历的现代政治发展，却在哲学中以观念的形式得到了反映。德国现实与德国哲学之间因此形成了一种明显的错位：一方面，现实社会仍然滞留在较为落后的政治状态之中；另一方面，德国哲学却已经在思想层面进入了现代国家所面临的问题。

因此，对德国而言，现实批判并不能简单照搬英法的路径。英法所要面对的现代国家问题，在德国尚未完全成为现实，但已经成为哲学问题。德国国家哲学和法哲学实际上构成了现代国家现实在德国思想领域中的理论反映，因此，对于德国而言，对这种其哲学的批判本身也就具有了现实批判的意义。

我们可以看到，马克思紧接着说：

\begin{adjustwidth}{2em}{2em}
\qquad\fangsong
“如果想从德国的现状[status quo]本身出发，即使采取唯一适当的方式，就是说采取否定的方式，结果依然是\textbf{时代错乱}。 即使对我国当代政治状况的否定，也已经是现代各国的历史废旧物品堆藏室中布满灰尘的史实。即使我否定了敷粉的发辫，我还是要同没有敷粉的发辫打交道。 即使我否定了1843年的德国制度，但是按照法国的纪年，我也不会处在1789 年，更不会是处在当代的焦点。”
\end{adjustwidth}

我们认为，马克思在这里想强调的是德国现实发展的严重滞后性。德国当时的社会现实，在英法等现代国家那里早已经成为历史。就是说，即使德国直接推翻当时落后的封建专制制度，也并不意味着德国就能够立即进入现代资本主义社会。德国所缺少的并不仅仅是一次政治制度上的否定，而是整个现代社会长期发展所形成的现实条件。换言之，德国当时通向现代资本主义社会的物质基础和社会条件尚不充分，单纯改变政治制度，并不能自动完成整个社会结构的现代转型。\footnote{编者注：在我们讨论的过程中，我们曾把当时德国的这种处境与中国清朝末年的情况作了一个有限的类比。清政府的统治即使被推翻，也并不意味着中国能够由此进入到资本主义社会（历史事实也是如此，清王朝的覆灭并没有意味着中国进入了资本主义社会），因为当时中国的社会经济结构、生产方式以及相应的社会条件并不会随着旧政权的倒台而自动发生根本转变。}

\vspace{1em}
此外，虽然马克思在前文用了很大篇幅说明德国现存制度的落后性，并认为这种制度已经远远落后于时代的发展，但马克思紧接着又明确提出要向德国的现实制度“开火”：
\begin{adjustwidth}{2em}{2em}
\qquad\fangsong
“向德国制度开火！一定要开火！这种制度虽然低于历史水平，低于任何批判，但依然是批判的对象，正像一个低于做人的水平的罪犯，依然是刽子手的对象一样。”
\end{adjustwidth}

这里就出现了一个很有意思的问题：既然德国制度已经落后于历史发展的水平，甚至已经“低于任何批判”，而且马克思也已经在前文中说了要首先对德国的哲学进行批判，那么为什么马克思在这里还要如此坚决地提出“向德国制度开火”？

经过我们的讨论，除去马克思个人的“写作情感”\footnote{编者注：我们在讨论中也认为，作为一名青年革命家，马克思可能写到这里写上头了。}，我们认为可能是由于下面的原因：

一个制度在历史上是否已经过时，与这个制度在现实中是否仍然存在，是两个不同的问题。德国的封建专制制度从整个欧洲历史发展的角度来看，固然已经是一种落后于时代的制度，甚至它所代表的历史原则在理论上也已经被现代社会的发展所否定；但是，这种制度在德国却仍然现实地存在着，并且仍然支配着德国社会。因此，不能因为它在历史上已经过时，就认为它不再值得批判。也就是说，即使首先要对德国哲学进行批判，但最后也是要落脚到德国现实的。因此，无论如何都有必要向德国制度“开火”。此外，作为一篇为后续批判工作奠定立场和方向的文章，马克思在这里也并不只是进行单纯的理论分析，也需要表明自己的基本态度。因此，“向德国制度开火！一定要开火！”这种强烈的表达，也可以从《导言》本身的文体特点来理解。

\vspace{1em}
之后，马克思在后文多处论述了德国现实的落后性，例如：
\begin{adjustwidth}{2em}{2em}
\qquad\fangsong
“这是一幅什么景象呵！ 社会无止境地继续分成各色人等，这些心胸狭窄、心地不良、粗鲁平庸之辈处于互相对立的状态，这些人正因为相互采取暧昧的猜疑的态度而被自己的统治者一律——虽然形式有所不同——视为 特予恩准的存在物。甚至他们还必须承认和首肯自己之被支配、被统治、被占有全是上天的恩准 ！”
\end{adjustwidth}

我们可以看到，在马克思看来，德国现实的落后并不只是表现为政治制度层面的落后，同时也表现为整个社会所呈现出来的落后状态。在封建专制制度以及等级分化的社会结构之下，社会被分割为彼此隔绝、彼此猜疑的不同群体，人与人之间缺乏真正的社会联合，而广大社会成员的主体意识和政治意识也受到严重束缚。

值得注意的是，马克思在这里指出，人们不仅事实上处于“被支配、被统治、被占有”的状态之中，甚至还要把这种状态理解为一种“上天的恩准”。这说明，落后的社会制度并不仅仅从外部支配着人，而且还会在一定程度上塑造人的意识，使人把现实中的支配关系当作一种自然的、合理的甚至不可改变的秩序。\footnote{编者注：毛主席曾说：“哪里有压迫，哪里就有反抗。”但在当时的德国，情况显然更加“微妙”。压迫的存在并不必然立即转化为自觉的反抗。当既有的统治关系被长期内化为一种理所当然的社会秩序时，被统治者甚至可能在观念上接受自身所处的地位，把现实中的不平等和支配关系理解为某种天然如此、不可违抗的安排。因此，德国革命的条件不仅在于人们确确实实在客观上受到压迫，还在于人们是否能够认识到这种压迫，并形成改变现实的主体意识。}

因此，马克思会在后文强调要让德国民众意识到自身处在被压迫的状况之下，意识到现存社会制度的落后性，从而激发自身的反抗意识：

\begin{adjustwidth}{2em}{2em}
\qquad\fangsong
“问题在于不让德国人有一时片刻去自欺欺人和俯首听命。应当让受现实压迫的人意识到压迫，从而使现实的压迫更加沉重；应当公开耻辱，从而使耻辱更加耻辱。应当把德国社会的每个领域作为德国社会的羞耻部分[partie honteuse]加以描述，应当对这些僵化了的关系唱一唱它们自己的曲调，迫使它们跳起舞来！为了激起人民的勇气，必须使他们对自己大吃一惊。这样才能实现德国人民的不可抗拒的要求，而各国人民的要求本身则是能使这些要求得到满足的决定性原因。”
\end{adjustwidth}

之后，马克思还进一步指出，对德国落后制度的批判和反抗，并不仅仅具有德国自身的意义，对于已经进入现代资本主义社会的国家同样具有启示作用。马克思是这么论述的：

\begin{adjustwidth}{2em}{2em}
\qquad\fangsong
“甚至对现代各国来说，这种反对德国现状的狭隘内容的斗争，也不会是没有意义的，因为德国现状是旧制度 [ancien régime] 的公开的完成，而旧制度是现代国家的隐蔽的缺陷。对当代德国政治状况作斗争就是对现代各国的过去作斗争，而对过去的回忆依然困扰着这些国家。这些国家如果看到，在它们那里经历过自己的悲剧的旧制度，现在又作为德国的幽灵在演自己的喜剧，那是很有教益的。”
\end{adjustwidth}

我们认为,马克思在这里的意思是，德国虽然在现实发展上落后于英法等现代国家，但正因为德国仍然以一种较为赤裸、直接的形式保留着旧制度，所以许多在现代国家中已经被新的政治形式所遮蔽的问题，反而能够在德国这里更加明显地表现出来。换言之，现代资本主义国家虽然已经完成了政治制度上的现代化，但这并不意味着旧社会中的一切矛盾都已经被真正消除，其中一些支配关系和社会矛盾只是以新的形式继续存在。因此，对德国现存制度的批判对其他现代国家而言也是具有极为重要的启示意义的。

\subsection{本期小结}
本期读书会是这学期的第一次读书会，也是咱们师门第一次尝试用讨论的形式来开展经典著作研读。整体来说，我觉得这次的氛围非常好，这一个半小时几乎没有冷场的时候，大家讨论得都非常积极，很多问题都是在一来一回的交流中慢慢展开的。我自己比较深的一个感受是，有些句子一个人读的时候觉得好像已经看懂了，但真正拿到桌面上讨论，才会发现里面其实还有很多值得继续追问的地方。大家对一些问题的理解并不完全一样，有时候甚至会出现不同意见，但这种不同反而让讨论变得更有意思。很多原本比较模糊的认识，也是在彼此补充和追问的过程中逐渐清晰起来的。这次能够负责本学期的读书会，对我来说也是一次很难得的锻炼机会。非常感谢刘老师和师门各位同学对我的信任，愿意把这项工作交给我，也感谢刘老师给了我比较大的空间，让我能够尝试一些新的形式。第一次组织这样的读书会，肯定还有很多不成熟的地方，但从这次的效果来看，我觉得这种形式还是很值得继续做下去的。也希望接下来的几期能够在这次的基础上继续改进，把经典读得更细一些，把问题讨论得更深入一些，也让每一次读书会都能够真正有所收获。最后，我觉得本次读书会还是收获满满的。也欢迎大家对之后读书会的内容和形式多提意见、多提建议，更欢迎大家继续参加接下来的读书会。希望我们能够一起把这个读书会越办越好！

\newpage

\section{第二期：《<黑格尔法哲学批判>导言》（2）}
\subsection{读书会记录}
2026年10月17日13：30—15：00，我们在同济大学马克思主义学院413教室开展了第二期的读书会。本期读书会我们对《<黑格尔法哲学批判>导言》上期剩下的部分进行了阅读，并展开了充分的讨论。具体讨论内容可概括如下：

\subsubsection{一、应怎样推动德国革命？已有方案存在着何种局限？}\label{sec:2.1}

承接第一期关于德国现实制度落后性的讨论，马克思在后文提出了一句十分重要的话：

\begin{adjustwidth}{2em}{2em}
\qquad\fangsong
“我们德国人幸而不是西徐亚人。”
\end{adjustwidth}

马克思在前文提到，希腊哲学家中虽然有一个是西徐亚人，但并不能由此认为西徐亚人整体上已经进入了希腊文化所达到的历史水平。这里要区分的是\textbf{个别人在思想文化上的成就}与\textbf{一个民族或国家的整体历史发展水平}。

那么，德国人为什么“幸而不是西徐亚人”？马克思紧接着作出了解释：

\begin{adjustwidth}{2em}{2em}
\qquad\fangsong
“正像古代各民族是在想象中、在神话中经历了自己的史前时期一样，我们德国人在思想中、在哲学中经历了自己的未来的历史。我们是当代的哲学同时代人，而不是当代的历史同时代人。德国的哲学是德国历史在观念上的延续。”
\end{adjustwidth}

我们可以看到，马克思在这里实际上延续了我们在第一期所讨论的德国发展的特殊性的观点：\textbf{德国的现实政治发展落后于英法等现代国家，但德国的哲学却已经在观念上思考了现代国家所面临的问题。}因此，德国不仅仅是偶然出现了几个思想先进的哲学家，而是在哲学领域形成了一套能够与现代国家的政治发展进行理论对话的思想体系。

用通俗一点的话说，英法等国在现实中经历的政治和社会变革，德国人却在很大程度上通过哲学思辨的方式提前“经历”了。德国社会在现实发展中尚未完成的某些历史任务，在哲学中却已经成为被思考和讨论的问题。因此，马克思才会说，德国人是当代的“哲学同时代人”，却不是当代的“历史同时代人”。

但是，这里又产生了一个问题：既然德国的哲学已经达到了现代国家的思想水平，那么是不是只要把这种哲学直接应用于德国社会，就能够实现德国的现代化？马克思显然并不这么认为。\textbf{哲学上的先进并不意味着现实社会的发展条件已经成熟，更不意味着哲学观念能够自动转化为现实制度。}正因如此，德国的出路既不能仅仅停留于批判落后的现实制度，也不能只是要求现实去实现既有的哲学观念，而必须进一步批判哲学本身，并考察这种哲学能够成为现实的条件。

紧接着，马克思讨论了德国当时两种不同的政治思想倾向：\textbf{实践政治派}和\textbf{起源于哲学的理论政治派}\footnote{编者注：根据马恩选集第一卷的注释，这里的实践政治派包括一部分自由资产阶级和知识分子以及民主派的代表，他们提出了实践政治的研究。这里的理论政治派带有整个青年黑格尔运动的特征，他们从黑格尔哲学得出彻底的无神论结论，但同时又使哲学脱离现实，从而在事实上日益脱离实际革命斗争。}。对于这两派，马克思都采取了批判的态度，但是批判的侧重点并不相同。

首先，马克思论述了实践政治派的错误：

\begin{adjustwidth}{2em}{2em}
\qquad\fangsong
“因此，德国的实践政治派要求对哲学的否定是正当的。该派的错误不在于提出了这个要求，而在于停留于这个要求——没有认真实现它，也不可能实现它。”
\end{adjustwidth}

在马克思看来，实践政治派强调现实的政治斗争、反对停留于哲学思辨，这一要求本身具有合理性。因为社会的解放终究不能只是在思想中完成，必须落实到现实的政治和社会生活中。但是，它的问题在于把对哲学的否定理解得过于简单，以为只要不再谈论哲学，或者对哲学表达一些不满，就已经真正超越了哲学。

马克思对此作了一个十分精辟的概括：

\begin{adjustwidth}{2em}{2em}
\qquad\fangsong
“一句话，你们不使哲学成为现实，就不能够消灭哲学。”
\end{adjustwidth}

这句话应当如何理解？我们认为，马克思这里所说的“消灭哲学”，并不能简单理解为不再研究哲学，更不是要否定一切哲学思想。联系上下文来看，马克思批判的主要是\textbf{把哲学作为一种脱离现实的、独立自足的思辨活动}。如果社会现实中仍然存在哲学所反映和思考的矛盾，那么无论我们在口头上怎样宣布哲学已经无用，都不能真正解决这些矛盾。只有把哲学中具有现实意义的解放要求转化为实际的社会变革，哲学作为与现实相分离的抽象要求才可能被真正扬弃。

随后，马克思又转向起源于哲学的理论政治派：

\begin{adjustwidth}{2em}{2em}
\qquad\fangsong
“该派认为目前的斗争只是哲学同德国世界的批判性斗争，它没有想到迄今为止的哲学本身就属于这个世界，而且是这个世界的补充，虽然只是观念的补充。”
\end{adjustwidth}

理论政治派的问题与实践政治派恰好相反。实践政治派看到了实践的重要性，却忽视了哲学与现实的联系；理论政治派则重视哲学的批判作用，却没有对哲学自身采取同样的批判态度。它把哲学作为批判现实的出发点，却没有进一步追问：\textbf{这种哲学本身是不是也产生于需要被批判的现实社会？其思想前提是不是同样存在局限？}

马克思因而把理论政治派的根本缺陷概括为：

\begin{adjustwidth}{2em}{2em}
\qquad\fangsong
“它以为，不消灭哲学，就能够使哲学成为现实。”
\end{adjustwidth}

由此我们能够看到两派错误之间的对称关系：\textbf{实践政治派企图不实现哲学就消灭哲学，理论政治派则企图不批判和扬弃哲学就实现哲学。}前者忽视了理论所包含的现实问题，后者忽视了理论自身的历史局限。马克思所提出的要求，并不是在“重视理论”与“重视实践”之间作简单选择，而是要求理论批判与现实变革相互贯通。


\subsubsection{二、马克思的观点：要从理论批判上升到实践批判}

在指出两种政治派别的局限之后，马克思开始明确自己的理论批判对象：

\begin{adjustwidth}{2em}{2em}
\qquad\fangsong
“德国的国家哲学和法哲学在黑格尔的著作中得到了最系统、最丰富和最终的表述；对这种哲学的批判既是对现代国家以及同它相联系的现实所作的批判性分析，又是对迄今为止的德国政治意识和法意识的整个形式的坚决否定……”
\end{adjustwidth}

为什么对一种哲学学说的批判，会同时成为对现代国家现实的批判？这看起来似乎与前面强调的现实批判又有些矛盾，但如果把前后文联系起来，问题就比较容易理解了。

首先，黑格尔的法哲学并不是纯粹凭空构造出来的观念体系，它包含着对现代国家、法、政治制度和市民社会等问题的哲学思考。从这个意义上说，\textbf{对黑格尔法哲学的批判并非只是在哲学概念内部展开的争论，而是通过批判这些概念，进一步考察它们所表达和反映的现代社会关系。}

其次，马克思并不认为黑格尔法哲学已经正确地解决了现代国家的问题。他指出，现代国家自身存在着一种与“现实的人”相分离的倾向，而德国关于现代国家的思辨哲学又把这种分离表现为抽象的思想形式。因此，对黑格尔法哲学的批判不仅要指出哲学论证中存在的问题，还要进一步追问：\textbf{为什么现代国家会以这样的方式被理解？这种哲学形式究竟反映了现代国家怎样的现实矛盾？}

这一点也可以帮助我们理解马克思之后所说的：

\begin{adjustwidth}{2em}{2em}
\qquad\fangsong
“对思辨的法哲学的批判既然是对德国迄今为止政治意识形式的坚决反抗，它就不会专注于自身，而会专注于课题，这种课题只有一个解决办法：实践。”
\end{adjustwidth}

这里的关键在于，\textbf{批判的出发点可以是哲学，但其最终指向不能停留在哲学本身。}理论批判的目的不是简单地用一个哲学体系取代另一个哲学体系，而是要揭示产生这些理论矛盾的现实基础，并把问题引向现实的社会实践。这样，马克思便从前面对德国哲学特殊地位的说明，进一步转向了理论与实践关系的问题。

之后，马克思提出了一段极为著名的论述：

\begin{adjustwidth}{2em}{2em}
\qquad\fangsong
“批判的武器当然不能代替武器的批判，物质力量只能用物质力量来摧毁；但是理论一经掌握群众，也会变成物质力量。”
\end{adjustwidth}

这句话乍一看很好理解：理论不能代替实践，只有依靠现实的物质力量才能改变既有的社会关系。但是，如果仅仅把它理解为“实践比理论更重要”，实际上又会忽视马克思接下来的论述。因为马克思紧接着又说，\textbf{理论一经掌握群众，也会变成物质力量。}

我们认为，理解这段话需要区分两个层面。

第一，理论批判与现实变革不是同一种活动。理论可以揭示一种制度为什么不合理，可以批判和否定既有制度的思想基础，却不能仅仅依靠思想上的批判就使现实存在的社会关系自动消失。例如，即使人们已经从理论上充分认识到一种社会制度的历史局限，这种制度所依赖的政治权力、经济关系以及社会组织仍然现实地存在着。因此，要改变这种制度，最终还需要能够作用于现实的社会力量。

第二，理论并不是与物质力量毫无关系。理论本身虽然不能直接代替现实行动，但当它能够揭示现实矛盾、表达社会成员的实际需要，并且为广大群众所理解和掌握时，就可能转化为指导、组织和推动现实行动的力量。\textbf{也就是说，理论转化为物质力量，并不是因为思想能够凭空创造现实，而是因为掌握理论的人能够在一定的社会条件下开展现实的实践活动。}

因此，马克思随后又说：

\begin{adjustwidth}{2em}{2em}
\qquad\fangsong
“理论只要说服人，就能掌握群众；而理论只要彻底，就能说服人。所谓彻底，就是抓住事物的根本。而人的根本就是人本身。”
\end{adjustwidth}

这里又出现了一个值得进一步思考的问题：为什么马克思要强调理论的“彻底性”？我们认为，这里的“彻底”并不是说理论的措辞越激烈越好，而是说理论必须能够超越表面的现象，把握问题产生的根本原因。\textbf{只有真正触及人的现实处境和现实解放需要的理论，才有可能获得广泛的认同，并通过群众的实践转化为现实力量。}

可见，理论的力量来源于它对现实问题的把握，也实现于群众改变现实的实践之中。马克思在这里把理论的批判性、群众的主体性和社会变革的实践性联系了起来。

\subsubsection{三、宗教改革与德国革命的关系}

在强调理论的实践作用之后，马克思又联系德国的历史经验，讨论了路德宗教改革的意义：

\begin{adjustwidth}{2em}{2em}
\qquad\fangsong
“即使从历史的观点来看，理论的解放对德国也有特殊的实践意义。德国的革命的过去就是理论性的，这就是宗教改革。正像当时的革命是从僧侣的头脑开始一样，现在的革命则从哲学家的头脑开始。”
\end{adjustwidth}

为什么要在这里提到宗教改革？我们认为，马克思并不是要把宗教改革与他所追求的现实解放完全等同起来，而是想借助德国历史上的一个重要事件说明：\textbf{在德国，思想领域的变革曾经成为推动社会历史变化的重要起点。}既然德国具有这样的历史特点，那么理论批判对于当时德国革命就具有不可忽视的意义。

不过，马克思并没有因此对路德的宗教改革作出全盘肯定。恰恰相反，他紧接着指出：

\begin{adjustwidth}{2em}{2em}
\qquad\fangsong
“路德战胜了虔信造成的奴役制，是因为他用信念造成的奴役制代替了它。他破除了对权威的信仰，是因为他恢复了信仰的权威。他把僧侣变成了世俗人，是因为他把世俗人变成了僧侣。”
\end{adjustwidth}

这段话是什么意思？简单来说，在马克思看来，路德的宗教改革虽然对传统教会的权威和旧有的宗教关系形成了冲击，却没有真正消除宗教对人的束缚，而是使这种束缚从外在的宗教权威转向了人的内在信仰。人们表面上似乎摆脱了外部权威的支配，但对宗教的依赖仍然以另外一种形式继续存在。

因此，宗教改革虽然具有历史进步意义，却并没有真正实现人的解放。马克思对此作了一个辩证的判断：“新教即使没有正确解决问题，毕竟正确地提出了问题。”

这也可以帮助我们理解马克思为什么在前面既强调理论的作用，又强调理论不能代替现实变革。\textbf{思想上的解放可以成为现实解放的起点，但如果思想上的变革没有进一步转化为现实社会关系的变革，那么人的解放仍然是不彻底的。}德国过去的宗教改革是如此，马克思所期待的德国革命同样必须面对这一问题。

\subsubsection{四、德国革命的现实性}\label{sec:2.6}

如果说前面的论述主要强调理论对于德国革命的重要意义，那么马克思接下来的分析，则进一步指出理论实现自身所面临的现实困难。他说：

\begin{adjustwidth}{2em}{2em}
\qquad\fangsong
“可是，彻底的德国革命看来面临着一个重大的困难。就是说，革命需要被动因素，需要物质基础。理论在一个国家实现的程度，总是取决于理论满足这个国家的需要的程度。”
\end{adjustwidth}

马克思为什么在强调理论的彻底性之后，又突然强调革命需要“物质基础”？我们认为，这实际上是在防止把前面对理论作用的肯定理解为一种唯心主义式的观点。\textbf{理论能否实现，并不单纯取决于理论本身是否正确、是否先进，还取决于现实社会是否存在使这种理论获得实践力量的条件。}

为了进一步说明这一点，马克思又提出：

\begin{adjustwidth}{2em}{2em}
\qquad\fangsong
“光是思想力求成为现实是不够的，现实本身应当力求趋向思想。”
\end{adjustwidth}

所谓“思想力求成为现实”，指人们在观念上提出某种社会变革的要求，并试图使这种要求得到实现。但马克思强调，仅有这样的主观愿望是不够的。所谓“现实本身应当力求趋向思想”，可以理解为\textbf{现实社会的发展必须产生与这种理论要求相联系的矛盾、需要以及能够推动变革的社会力量。}

这里举一个比较容易理解的例子。即使一个人提出了某种非常先进的社会理想，但如果当时社会的生产方式、阶级关系以及现实需要完全不具备实现这种理想的条件，那么这种理想就很难仅凭思想家的努力成为现实。反过来说，当现实社会已经产生迫切需要解决的问题，而一种理论又能够准确把握这种需要，并为实际行动提供方向时，理论便有可能获得实践基础。

当然，我们也不能把马克思在这里的表述简单理解为“只有等到一切客观条件完全成熟之后才可以行动”。马克思强调的是理论与现实条件之间的联系，而不是否定人的能动性。理论能够在现实矛盾的基础上发挥作用，但它不能不顾历史条件而任意创造自己所需要的社会基础。


紧接着，马克思在讨论德国革命的可能性时，提出了一个看起来有些反常识的判断：

\begin{adjustwidth}{2em}{2em}
\qquad\fangsong
“对德国来说，彻底的革命、普遍的人的解放，不是乌托邦式的梦想，相反，局部的纯政治的革命，毫不触犯大厦支柱的革命，才是乌托邦式的梦想。”
\end{adjustwidth}

一般来说，我们可能会认为，局部的政治变革所涉及的范围更小，因而更容易实现；彻底的社会革命涉及的范围更大，因而实现起来更困难。但马克思为什么在这里作出了相反的判断？

理解这个问题，首先要看马克思如何解释“局部的纯政治的革命”的社会基础。在他的分析中，这样的革命通常意味着市民社会中的一个特殊阶级实现自身的解放，并以自己的特殊利益和政治要求代表整个社会的普遍利益。例如，法国资产阶级在反对封建贵族和僧侣特权的斗争中，曾经能够以某种普遍解放的姿态获得其他社会群体的支持。

但是，一个特殊阶级要想承担这样的历史角色，并不是只要宣布自己代表全社会就够了。它至少需要使自身的要求在特定历史条件下与广大社会成员的要求联系起来，并且在现实斗争中形成较为广泛的社会力量。因此，马克思会这么写道：

\begin{adjustwidth}{2em}{2em}
\qquad\fangsong
“只有为了社会的普遍权利，特殊阶级才能要求普遍统治。”
\end{adjustwidth}

那么，德国当时具不具备这样的条件？马克思的回答是否定的。他认为，当时德国市民社会中的各个阶级缺乏足够的革命精神和广泛的社会代表性，尤其是德国中等阶级并不能像法国资产阶级曾经做过的那样，把自己的特殊要求上升为能够广泛动员社会的普遍要求。德国的阶级关系也更加复杂：一些阶级还没有充分完成对旧势力的斗争，就已经面临来自其他阶级的新的对立。

这也就意味着，德国很难简单重复法国曾经经历过的、由某个特殊阶级领导并以政治解放为主要内容的历史道路。\textbf{对德国而言，旧制度的特殊矛盾已经同现代社会更普遍的矛盾交织在一起，因而仅仅通过某个阶级的局部解放，未必能够解决整个社会的解放问题。}

不过，这里也需要注意，马克思作出上述判断是针对19世纪40年代德国特定的政治经济状况而言的，不能把它直接上升为任何国家在任何时期都只能通过同一种革命道路实现社会解放的普遍公式。我们需要把马克思的论述放在当时德国与法国不同的社会发展条件下理解。


在论证德国的特殊阶级难以通过局部政治革命实现普遍解放之后，马克思进一步提出了一个决定性的追问：

\begin{adjustwidth}{2em}{2em}
\qquad\fangsong
“那么，德国解放的实际可能性到底在哪里呢？”
\end{adjustwidth}

对此，马克思的回答是：德国需要形成一个由于遭受普遍苦难而具有普遍性质的社会领域，一个只有通过解放整个社会才能获得自身解放的阶级。他最终明确指出：

\begin{adjustwidth}{2em}{2em}
\qquad\fangsong
“社会解体的这个结果，就是无产阶级这个特殊等级。”
\end{adjustwidth}

为什么马克思要把无产阶级作为德国实现普遍的人的解放的现实力量？我们认为，这里的关键在于马克思关于无产阶级社会地位的分析。无产阶级与其他特殊阶级不同的地方，在于它所遭受的并不是某种可以通过获得局部政治特权就彻底解决的特殊不公正。在马克思的分析中，无产阶级的形成直接关系到整个既有社会制度的矛盾，因而它的解放不能仅仅表现为在现存制度内部取得更有利的地位。

马克思用一段带有强烈哲学意味的话描述这一点：

\begin{adjustwidth}{2em}{2em}
\qquad\fangsong
“形成一个并非市民社会阶级的市民社会阶级，形成一个表明一切等级解体的等级，形成一个由于自己遭受普遍苦难而具有普遍性质的领域……”
\end{adjustwidth}

乍一看，“并非市民社会阶级的市民社会阶级”似乎是一种自相矛盾的说法。但如果联系上下文，就会发现它体现了马克思对无产阶级特殊历史处境的理解：无产阶级当然是在市民社会内部形成的阶级，但它的存在又集中表现了这个社会无法在自身既有关系中解决的矛盾。因此，它并不能像某些特殊阶级那样，仅仅通过占据现存社会中的一个更高位置来完成自身的解放。

在这个意义上，\textbf{无产阶级的解放要求具有超出自身特殊利益的普遍指向：它只有改变使自身成为无产阶级的社会关系，才能真正实现自身的解放。}马克思由此把无产阶级的自我解放与更广泛的人的解放联系起来。

当然，我们也需要注意，马克思在《导言》中对于无产阶级的认识仍然具有鲜明的思想形成时期的特点。这里主要是从普遍苦难、人的丧失以及普遍解放的角度来论证无产阶级的历史意义；至于无产阶级在资本主义生产关系中的具体地位，以及它如何在现实的阶级斗争中组织和实现自身的解放，则需要联系马克思之后的政治经济学研究和历史分析进一步理解。\textbf{不能把《导言》中的哲学性论证直接等同于马克思后来已经成熟的全部阶级理论。}


\subsubsection{五、为什么说“解放的头脑是哲学，心脏是无产阶级”？}

在《导言》的结尾，马克思把哲学批判和无产阶级的历史作用联系起来，提出：

\begin{adjustwidth}{2em}{2em}
\qquad\fangsong
“德国人的解放就是人的解放。这个解放的头脑是哲学，它的心脏是无产阶级。哲学不消灭无产阶级，就不能成为现实；无产阶级不把哲学变成现实，就不可能消灭自身。”
\end{adjustwidth}

为什么说哲学是“头脑”，无产阶级是“心脏”？我们认为，马克思在这里实际上把前文的两条论证线索结合在了一起。一方面，哲学提出了关于人的解放的理论要求，可以在思想上揭示旧制度及现代社会的矛盾，为现实变革提供理论方向。另一方面，单纯的哲学批判无法自动完成社会革命；现实的解放需要具有改变既有社会关系的利益要求和实践能力的社会力量。在马克思的论述中，无产阶级正是这样一种现实力量。

因此，哲学与无产阶级之间并不是一种单向的关系，即不是哲学家先构造好一个完整的方案，然后由无产阶级去机械地加以执行；也不能理解为无产阶级只要凭借自发的现实行动就必然能够完成普遍的人的解放。\textbf{马克思强调的是解放理论与现实社会力量之间的相互联系：理论需要在现实实践中得到实现，社会力量也需要在对自身处境和解放任务的认识中展开行动。}

这里还有一个容易引起误解的问题：马克思为什么要说“哲学不消灭无产阶级，就不能成为现实”？这里的“消灭无产阶级”是不是指消灭作为具体个人的无产者？显然不能这样理解。联系上下文可以看到，马克思讨论的是无产阶级作为特定社会关系产物的历史地位。因而，这里的“消灭无产阶级”，应当理解为\textbf{通过改变使人们处于无产阶级地位的社会关系，扬弃无产阶级作为一个受压迫阶级的存在条件}，而不是对无产阶级成员个人的否定。

同样，“无产阶级不把哲学变成现实，就不可能消灭自身”也不是说无产阶级需要简单接受哲学家的一切主张，而是说无产阶级要实现自身解放，就不能只满足于对自身处境的局部改善，还需要把人的解放这一普遍要求落实到改变现实社会关系的实践中。

至此，我们可以把《导言》后半部分的论证串联起来：\textbf{德国现实的落后性使哲学批判具有特殊意义；哲学批判只有指向现实实践才可能实现自身；现实革命既需要理论，也需要物质基础；而无产阶级则被马克思视为能够把普遍解放的理论要求与实际社会变革联系起来的历史力量。}

回过头再看第一期开始时马克思所说的“对宗教的批判是其他一切批判的前提”，就会发现《导言》的论证并不是停留于宗教问题本身。马克思实际上沿着\textbf{宗教批判—法和政治的批判—思辨哲学的批判—革命实践—无产阶级与人的解放}这样一条线索，逐步把问题从观念领域引向现实社会及其变革的可能性。

\subsection{本期小结}

第二期读书会延续了第一期的讨论，也把我们对《导言》的理解进一步引向了理论与实践、政治解放与人的解放以及无产阶级的历史地位等问题。与前半部分相比，后半部分的论述更加集中，马克思提出的许多判断也更加具有哲学上的概括性。经过前后两期的阅读，我觉得一个比较重要的收获是，我们不能仅仅把《导言》当作一篇包含许多著名论断的文章来读，而需要努力把握马克思的整个论证过程。马克思为什么从宗教批判谈起，为什么转而讨论德国的法哲学，为什么要批判当时的实践政治派和理论政治派，又为什么最终把问题落到无产阶级和人的解放上？这些问题并不是彼此孤立的，而是构成了一条逐步展开的思想线索。只有把这些线索联系起来，我们才能较为完整地理解这篇著作的理论意义。

同时，我们也需要注意到，《导言》写于马克思思想发展的重要转折时期。文章已经非常鲜明地表现出马克思对现实社会关系、革命实践和人的解放问题的关注，但其中关于无产阶级、革命和人的本质等问题的论证仍然需要放在马克思思想发展的历史过程中加以理解。我们不应简单地用马克思后期的成熟理论去替代这篇文章自身的思想内容，也不应把其中的哲学判断直接视为不需要进一步论证的最终结论。

至此，我们完成了《<黑格尔法哲学批判>导言》的两期研读。对我来说，读书会最有价值的地方，仍然在于通过共同阅读、反复追问和相互讨论，把原本看似已经理解的内容重新变成值得深入思考的问题。希望之后的经典著作研讨也能够继续坚持这种方式：不仅重视结论，更重视结论是怎样得到的；不仅重视对原文的解释，也重视对我们自己理解的反思。总之，本期读书会也是收获满满的，期待在接下来的读书会中继续和大家一起讨论。

\newpage
\end{document}


\end{document}
